\chapter{Algebraic theories and e-graphs}\label{chapter:algebras}

In this chapter we will present an algebraic view of e-graphs, which will be the basis for our generalisation in later chapters.
The chapter introduces the notion of an algebraic theory and recalls e-graphs as a data structure for representing congruences on terms.
\section{Algebraic theories}
Traditionally, these are defined based on a signature $\Sigma$.
\begin{definition}
An (algebraic) signature $\Sigma$ is given by a function $\Sigma : \mathbb{N} \to \catname{Set}$.
\end{definition}

Intuitively, $\Sigma(n)$ returns the set of generators of arity $n$.
Equivalently, it can be thought of as a set of generators with specified arity, such that for all arities with no representative in the set $\Sigma(n) = \varnothing$.

\begin{definition}\label{def:sigma-alg}
Let $\Sigma$ be a signature.
$\Sigma$-\emph{algebra} is a set $A$ together with a family of operations indexed by $n \in \mathbb{N}$
\[
\alpha_n : \Sigma(n) \times A^n \to A .
\]
A homomorphism of $\Sigma$-algebras is a map $f : A \to B$ such that for each $n \in \mathbb{N}$, the following diagram commutes
% https://q.uiver.app/#q=WzAsNCxbMCwwLCJcXFNpZ21hKG4pXFx0aW1lcyBBXm4iXSxbMiwwLCJBIl0sWzAsMiwiXFxTaWdtYShuKVxcdGltZXMgQl5uIl0sWzIsMiwiQiJdLFswLDEsIlxcYWxwaGFfbiJdLFsxLDMsImYiXSxbMiwzLCJcXGJldGFfbiIsMl0sWzAsMiwiXFxpZCBcXHRpbWVzIGZebiIsMl1d
\[\begin{tikzcd}
	{\Sigma(n)\times A^n} && A \\
	\\
	{\Sigma(n)\times B^n} && B
	\arrow["{\alpha_n}", from=1-1, to=1-3]
	\arrow["{\id \times f^n}"', from=1-1, to=3-1]
	\arrow["f", from=1-3, to=3-3]
	\arrow["{\beta_n}"', from=3-1, to=3-3]
\end{tikzcd}\]
$\Sigma$-algebras and their homomorphisms form a category $\catname{\Sigma}\text{-}\catname{Alg}$.
\end{definition}
Put simply, $\alpha_n$ takes an $n$-ary generator and interprets it in the set $A$: partially applying $\alpha_n$ to $\Sigma(n)$ yields a function $A^n \to A$, i.e., an interpretation.
\begin{example}\label{example:sigma_algebra}
Many familiar constructions are instances of $\Sigma$-algebras.
\begin{itemize}
\item A monoid is a $\Sigma$-algebra with $\Sigma(0) = \{e\}$, $\Sigma(2) = \{\cdot\}$, and $\Sigma(n) = \varnothing$ otherwise.
\item A group is a $\Sigma$-algebra with $\Sigma(0) = \{e\}$, $\Sigma(1) = \{^{-1}\}$, $\Sigma(2) = \{\cdot\}$, and $\Sigma(n) = \varnothing$ otherwise.
\item The \cref{example:term_rewrite} can be seen as a $\Sigma$-algebra with $\Sigma(0) = \{a\} \cup \mathbb{N}$, $\Sigma(2) = \{*, /, <\!\!<\}$, and $\Sigma(n) = \varnothing$ otherwise. 
\item In general, a programming language is not an instance of a $\Sigma$-algebra, as the latter, in particular, does not account for binding.
\end{itemize}
\end{example}
Then, for a given signature $\Sigma$ there exists a \emph{free} $\Sigma$-algebra which we define next.
\begin{definition}\label{def:free_sigma_algebra}
 Let $X$ be a set and $\Sigma$ be a signature.
 Let $T_{\Sigma}(X)$ be the smallest set generated by the rules below:
 \[
 \inferrule*[]{x \in X}{x \in T_{\Sigma}(X)} \qquad \inferrule*[]{f \in \Sigma(m) \quad t_0 \in T_{\Sigma}(X), \ldots, t_{m-1} \in T_{\Sigma}(X)}{f(t_0, \ldots, t_{m-1}) \in T_{\Sigma}(X)}~.
 \] 
 Then $(T_{\Sigma}(X), \tau_n)$ is a $\Sigma$-algebra, where $\tau_n : \Sigma(n) \times T_{\Sigma}(X)^n \to T_{\Sigma}(X)$ is given by the second rule above.
\end{definition}

\begin{proposition}
The set $T_{\Sigma}(X)$ is the free $\Sigma$-algebra over $X$.
\end{proposition}
Essentially, $T_{\Sigma}(X)$ is the set of all terms generated by the signature $\Sigma$ and the variables in $X$.
\begin{proposition}
The functor $T_{\Sigma}(-) : \catname{Set} \to \catname{\Sigma}\text{-}\catname{Alg}$ has a right adjoint 
\begin{gather*}
U : \catname{\Sigma}\text{-}\catname{Alg} \to \catname{Set}\\
(A, \alpha) \mapsto A
\end{gather*}
\end{proposition}
Using the properties of freeness, for a $\Sigma$-algebra $(A, \alpha)$ and a set $X$, we can define a map
\[
\hat{\alpha_{X}} : T_{\Sigma}(X) \times A^X \to A
\]
which \emph{evaluates} a term in $T_{\Sigma}(X)$ using the assignment of variables from $X$ in $A$, where
$\hat{\alpha_{X}}(-, a) : T_{\Sigma}(X) \to A$ is the morphism, adjoint to the variable assignment morphism $a : X \to A$.
Of course, the terms are only part of the picture, and, for example, all our example in \cref{example:sigma_algebra} come with certain equations, in addition to the terms. 
\begin{definition}
Let $\Sigma$ be a signature.
An equation in context $n$ is a pair $(t_1, t_2) \in T_{\Sigma}(n) \times T_{\Sigma}(n)$ of terms with the same context. 
\end{definition}
\begin{remark}
  As terms use only finitely many variables, it justifies writing $n$ instead of arbitrary $X$ in $T_{\Sigma}(n)$.
\end{remark}
\begin{definition}
Let $\Sigma$ be a signature.
A $\Sigma$-algebra $(A, \alpha)$ satisfies an equation $(t_1, t_2) \in T_{\Sigma}(n) \times T_{\Sigma}(n)$ if for all valuations $a : n \to A$, the following holds:
\[
\hat{\alpha_{n}}(t_1, a) = \hat{\alpha_{n}}(t_2, a) ,
\]
that is, $t_1$ and $t_2$ agree under all valuations in $A$.
\end{definition}
\begin{remark}
We can introduce a syntactic sugar
\[
\llbracket t \rrbracket \defeq \hat{\alpha_{n}}(t, -) : A^n \to A~.
\]
With this in hand $\llbracket t_1 \rrbracket = \llbracket t_2 \rrbracket$ means they agree as maps $A^n \to A$.
\end{remark}
\begin{proposition}\label{prop:homomorphism_commutes}
Let $(A, \alpha)$ and $(B, \beta)$ be $\Sigma$-algebras and $f : A \to B$ be a homomorphism.
Then $\llbracket - \rrbracket$ commutes with $f$, that is
\[
f \circ \llbracket t \rrbracket_A = \llbracket t \rrbracket_{B} \circ f^n
\]
\end{proposition}
\begin{definition}
  A (presentation of) algebraic theory is a pair $(\Sigma, \mathcal{E})$ given by a signature and an indexed family of equations in context $n$: $\mathcal{E}_{n} \subseteq T_{\Sigma}(n) \times T_{\Sigma}(n)$.
  The category $(\Sigma, \mathcal{E})\text{-}\catname{Alg}$ if the full subcategory of $\catname{\Sigma}\text{-}\catname{Alg}$ on those $\Sigma$-algebras that satisfy equations in $\mathcal{E}$.
\end{definition}

\begin{definition}\label{def:congruence}
Let $\Sigma$ be a signature and let $(A,\alpha)$ be a $\Sigma$-algebra.
A \emph{congruence} on $(A,\alpha)$ is an equivalence relation $R \subseteq A \times A$ that is closed under operations in $\Sigma(n)$, that is if
$(a_0, b_0), \ldots, (a_{n-1}, b_{n-1}) \in R$ and $f \in \Sigma(n)$, then $(f (a_0, \ldots, a_{n-1}), f(b_0, \ldots, b_{n-1})) \in R$.
\end{definition}
\begin{proposition}\label{prop:free_algebraic_theory}
The inclusion $(\Sigma, \mathcal{E})\text{-}\catname{Alg} \hookrightarrow \catname{\Sigma}\text{-}\catname{Alg}$ has a left adjoint.
\end{proposition}
\begin{proof}
Let $(A, \alpha)$ be a $\Sigma$-algebra, and let $R \subseteq A \times A$ be given by
\[
R = \{(\hat{\alpha}_{n}(t_1, a), \hat{\alpha}_{n}(t_2, a)) \;\vert\; (t_1, t_2) \in \mathcal{E}_n, a : n \to A\} .
\]
That is, $R$ is the set of pairs of elements in $A$ that need to be equal in order for $(A, \alpha)$ to satisfy the equations in $\mathcal{E}$.
Let $\bar{R}$ be the smallest congruence containing $R$.
Technically, it is built iteratively starting from $R$ by simultaneously closing $\bar{R}$ under reflexivity, symmetry, transitivity, and for $f \in \Sigma(n)$ and $(a_0, b_0), \ldots, (a_{n-1}, b_{n-1}) \in \bar{R}$ letting $(\alpha_{n}(f,a_0, \ldots, a_{n-1}), \alpha_{n}(f,b_0, \ldots, b_{n-1})) \in \bar{R}$.
We then define the quotient $\Sigma$-algebra $(A / \bar{R}, \alpha / \bar{R})$, where $\alpha / \bar{R}$ is defined by
\[
\alpha / \bar{R}(f, [a_0], \ldots, [a_{n-1}]) = [\alpha_{n}(f, a_0, \ldots, a_{n-1})]
\]
for all $f \in \Sigma(n)$ and equivalence classes $[a_i]$ in $A / \bar{R}$.
We then have that $\Sigma$-algebra $(A / \bar{R}, \alpha / \bar{R})$ satisfies the equations in $\mathcal{E}$.
The quotient map induces a homomorphism $(A, \alpha) \to (A / \bar{R}, \alpha / \bar{R})$ which is the unit of the adjunction.
We will call the functor $(-) / \bar{R} : (A, \alpha) \to (A / \bar{R}, \alpha / \bar{R})$ a \emph{reflector}.
\end{proof}
When $\Sigma$-algebra $(A, \alpha)$ is $(T_{\Sigma}(X), \tau)$, the quotient $(T_{\Sigma}(X) / \bar{R}, \tau / \bar{R})$ is the free $(\Sigma, \mathcal{E})$-algebra over $X$---obtained by combining the two adjunctions---that satisfies the equations in $\mathcal{E}$.
\begin{example}
Continuing \cref{example:sigma_algebra}, 
\begin{itemize}
\item a monoid is a $\Sigma$-algebra that satisfies the equations
\[
x \vdash x \cdot e \equiv x \equiv e \cdot x \qquad x, y, z \vdash (x \cdot y) \cdot z \equiv x \cdot (y \cdot z);
\]
\item a group is a $\Sigma$-algebra that satisfies the equations
\[
x \vdash x \cdot e \equiv x \equiv e \cdot x \qquad x, y, z \vdash (x \cdot y) \cdot z \equiv x \cdot (y \cdot z) \qquad x \vdash x^{-1} \cdot x \equiv e \equiv x \cdot x^{-1}
\]
\item the $\Sigma$-algebra of \cref{example:term_rewrite} satisfies the equations
\[
x \vdash x * 2 \equiv x <\!\!< 1 \qquad x \vdash x / x \equiv 1 \qquad x, y, z \vdash (x * y) / z \equiv x * (y / z) .
\]
\end{itemize}
\end{example}

Having constructed the congruence $\bar{R}$ in the proof of \cref{prop:free_algebraic_theory}, we next show how e-graphs compute (a part of) the reflector in the case of $(\Sigma, \mathcal{E})$-algebra $T_\Sigma(\varnothing)$, with equality saturation playing the role of the closure operation. 
We recall the definition of an e-graph due to Willsey et al.~\cite{EggPaper}.

Fix a signature $\Sigma$. The data of an e-graph is built from three layers.

\begin{definition}[E-nodes, e-classes, e-graph]
\label{def:egraph}
Let e-class \emph{identifiers} $a, b, \dots$ be drawn from some set of
opaque tokens. An \emph{e-node} is an operation symbol applied to a
tuple of e-class identifiers,
\[
n \;\Coloneqq\; f(a_1, \dots, a_m), \qquad f \in \Sigma(m),
\]
where the nullary case $m = 0$ gives a bare constant $f$. An
\emph{e-class} is a non-empty set of e-nodes, regarded as mutually
equal. An \emph{e-graph} $\egraph$ is a finite set of e-classes
together with an equivalence relation $\equivid$ on identifiers,
identifying the classes its e-nodes refer to.
\end{definition}

\noindent
The essential point of~\cref{def:egraph} is that an e-node
records its children as \emph{e-class identifiers}, not as e-nodes:
sub-expressions are referred to only up to the equivalence already
known, which is what lets one class stand for many terms and what makes
the structure compact. 
One should note, that by the very definition, an e-graph does not distinguish a set of variables and therefore \emph{represents} only closed terms.
We understand terms as inhabitants of $T_\Sigma(\varnothing)$ and below we define what it means for an e-graph to represent a term.

\begin{definition}[Represented terms]
\label{def:represented}
Representation is defined by recursion. An e-node $f(a_1, \dots, a_m)$
\emph{represents} a term $f(t_1, \dots, t_m)$ when each identifier
$a_i$ refers to an e-class that represents $t_i$. An e-class represents
a term if one of its e-nodes does, and the e-graph represents a term if
one of its e-classes does.
\end{definition}

\noindent
Representation induces an equivalence relation on terms, which is the
one of interest here: write $t_1 \equivterm t_2$ when $t_1$ and $t_2$ are
represented in the \emph{same} e-class. So an e-graph carries a set of
terms (those it represents) together with the relation $\equivterm$
identifying those it has been told, or has deduced, to be equal.

\begin{definition}[Congruence invariant]
\label{def:cong-invariant}
An e-graph satisfies the \emph{congruence invariant} when, for every
$f \in \Sigma(n)$ and identifiers $a_1, \dots, a_n$ and $b_1, \dots,
b_n$ with $a_i \equivid b_i$ for all $i$, the e-nodes $f(a_1, \dots, a_n)$
and $f(b_1, \dots, b_n)$ belong to the same e-class.
\end{definition}

\noindent
Equivalently: equal arguments yield equal applications, $a_i \equivid b_i
\Rightarrow f(a_1, \ldots, a_n) \equivid f(b_1, \ldots, b_n)$. An e-graph maintaining this
invariant represents not merely an equivalence relation on terms but a
\emph{congruence}---which is exactly the structure
~\cref{def:congruence} requires.

\begin{remark}[Implementation]
\label{rem:egraph-impl}
The definition above is the semantic skeleton; the realisation
of~\cite{EggPaper} adds the state needed to maintain it efficiently. An
e-graph there is a triple $(U, M, H)$: a union-find $U$ storing
$\equivid$, an e-class map $M$ sending each identifier to its e-class
(with $a \equivid b$ iff $M[a] = M[b]$), and a hashcons $H$ mapping
e-nodes to identifiers. A $\find$ operation canonicalises identifiers
to a representative, and an e-node is canonical when its children are;
the hashcons invariant requires $H$ to index every canonical e-node by
its class, so that an e-node congruent to a given one can be located in
near-constant time. None of this extra state changes which congruence
is represented---it is what makes the operations of
 efficient---so we suppress it in
the main development and refer to~\cite{EggPaper} for details.
\end{remark}

Given a term $t \in T_\Sigma(\varnothing)$ and a set of equations
$\mathcal{E}$, denote $A_0$ the set of all sub-terms of $t$.
For $T_\Sigma(\varnothing)$ to satisfy $(t_1, t_2) \in \mathcal{E}_n$, it
means that $\hat{\tau}_n(t_1, a) = \hat{\tau}_n(t_2, a)$ for all
valuations $a : n \to T_\Sigma(\varnothing)$.
Restricting to $A_0$, it suffices that the equalities hold for
valuations $a : n \to A_0$ into the represented terms.
This is what an e-graph stipulates: given the set of terms $A_0$, for
each $(t_1, t_2) \in \mathcal{E}_n$ and each valuation $a : n \to A_0$,
it forms the instances $\hat{\tau}_n(t_1, a)$ and $\hat{\tau}_n(t_2, a)$
and identifies them.

This single pass, however, does not compute the congruence we are after,
for two reasons. First, identifying the instance-pairs yields only an
equivalence relation; for it to respect the operations---so that
$f(\vec u\,)$ and $f(\vec v\,)$ are identified whenever their arguments
already are---the e-graph additionally maintains the \emph{congruence
invariant} of \cref{def:cong-invariant}, closing the relation under the
operations present. Second, and less obviously, the set of represented
terms does not stay fixed: the instance $\hat{\tau}_n(t_2, a)$ produced
by firing $(t_1, t_2)$ need not lie in $A_0$, and once added it---together
with its sub-terms---enlarges the set against which subsequent valuations
are taken. An equation can only be instantiated at terms already
present, and firing is what makes new terms present.

The e-graph therefore computes not a single relation on $A_0$ but an
increasing chain
\[
A_0 \subseteq A_1 \subseteq A_2 \subseteq \cdots,
\qquad
R_0 \subseteq R_1 \subseteq R_2 \subseteq \cdots,
\]
where at each stage it fires every available instance and then closes
under the congruence invariant. In doing so it departs from the literal
reading above in one further respect: a variable is instantiated not by
a represented term but by an \emph{e-class}, so the valuations range over
the current quotient,
\[
\sigma : n \longrightarrow A_k / R_k ,
\]
and firing forms the instance-classes $\overline{\tau}_n(t_1, \sigma)$
and $\overline{\tau}_n(t_2, \sigma)$, adding the right-hand side to pass
from $A_k$ to $A_{k+1}$ when its term is not yet represented, and adding
their pair to the relation. Both components are monotone. When the chain
reaches a fixed point, $A_{k+1} = A_k$ and $R_{k+1} = R_k$---the e-graph
is \emph{saturated}---no instance and no congruence step adds anything
further, so $R_k$ is closed under precisely the clauses defining $\bar R$
in the proof of \cref{prop:free_algebraic_theory}. Hence
\[
R_k \;=\; \bar R \!\restriction_{A_k},
\]
the restriction to the represented terms of the congruence generated by
$\mathcal{E}$ on $T_\Sigma(\varnothing)$.

Instantiating variables by e-classes rather than by terms loses nothing,
and is the reason the congruence invariant is worth maintaining. Because
$R_k$ is a congruence, evaluation factors through the quotient
$q_k : A_k \to A_k/R_k$: if $a \equiv_{R_k} a'$ then $\hat{\tau}_n(t_i,
a) \equiv_{R_k} \hat{\tau}_n(t_i, a')$, so $\hat{\tau}_n(t_i, -)$
descends to a map $\overline{\tau}_n(t_i, -) : (A_k/R_k)^n \to A_k/R_k$
with $q_k \circ \hat{\tau}_n(t_i, -) = \overline{\tau}_n(t_i, -) \circ
q_k^{\,n}$. A single class-valuation $\sigma$ thus stands for all the
term-valuations that agree up to the congruence so far---of which there
may be unboundedly many---and the e-graph instantiates each equation
once per class-valuation rather than once per term-valuation. Were $R_k$
a bare equivalence, evaluation would not descend to the quotient and
matching against e-classes would be unsound.

In the language of \cref{prop:free_algebraic_theory}, a saturated e-graph
computes the unit of the reflection
\[
q \;:\; (T_\Sigma(\varnothing), \tau)
\longrightarrow (T_\Sigma(\varnothing)/\bar R,\ \tau/\bar R),
\qquad t \mapsto [t],
\]
restricted to the finite set $A_k$ of terms it represents. Each
represented term $t$ is assigned its e-class, a finite representative of
the class $[t] = q(t)$, and two represented terms share an e-class if and
only if they have the same image under $q$. The e-classes are thus in
bijection with $q(A_k)$, a finite subset of the carrier of the initial
$(\Sigma, \mathcal{E})$-algebra $T_\Sigma(\varnothing)/\bar R$: the
fragment spanned by $t$ and the terms reachable from it under
$\mathcal{E}$. The e-graph does not hold this algebra, but computes a
finite window onto it.

Above we have seen a very abstract account of e-graph, but of course, in real applications the efficiency of computing the reflector is important.
We will review the key elements that make e-graphs efficient in this regard next.
\begin{example}
\label{ex:egraph-run}
We illustrate the three mechanisms at work---multi-site matching, the
addition of instantiated right-hand sides, and congruence
closure---on a single small run.
We also illustrate the action of union-find and hashcons components of the implementation, which are the key to efficiency.
\paragraph{Union-find.} The union-find $U$ maintains the equivalence $\equivid$ on identifiers by the means of the data structure introduced by Tarjan~\cite{Tarjan}.
The crux of the data structure are two operations: $\union$ and $\find$.
Given two identifiers $a$ and $b$, $\union(a,b)$ sets $a$ to be the parent of $b$.
Given an identifier $a$, $\find(a)$ returns the root, i.e., the canonical identifier, of the induced tree of identifiers, which makes it easy to check if two identifiers belong to the same e-class by checking if $\find(a) = \find(b)$.
\paragraph{Hashcons.} The hashcons, attributed to hashing for the \texttt{cons} operator in LISP~\cite{hashconsing}, is a map $H$ from \emph{canonical} e-nodes---those all of whose child identifiers are canonical---to the identifier of the e-class containing them.

Take the signature with a constant
$2$, a binary $\ast$, a binary $+$, and a unary $d$, and the single
equation
\[
x \ast 2 \;=\; x + x \qquad (x \in T_\Sigma(1)).
\]
As the seed take the closed term
\[
t \;=\; \bigl(d(2 \ast 2) + d(2 + 2)\bigr) \;+\; \bigl((2 \ast 2) \ast 2\bigr),
\]
whose sub-term closure $A_0$ consists of
\[
2,\quad 2\ast 2,\quad 2+2,\quad (2\ast 2)\ast 2,\quad d(2\ast 2),\quad
d(2+2),
\]
together with the two $+$-terms that assemble them into $t$.
Each sub-term corresponds to an e-node which occupies its own e-class
initially; we summarise the e-nodes and e-classes below, in the order
they are built. We write $\alpha_i$ for e-class identifiers.

\begin{center}
\begin{tabular}{c|c|c}
Term & E-node & E-class\\
\hline
2 & 2 & $\alpha_0$\\
$2 \ast 2$ & $\ast(\alpha_0, \alpha_0)$ & $\alpha_1$\\
$d(2 \ast 2)$    & $d(\alpha_1)$        & $\alpha_2$ \\
$2 + 2$          & $+(\alpha_0, \alpha_0)$    & $\alpha_3$ \\
$d(2 + 2)$       & $d(\alpha_3)$              & $\alpha_4$ \\
$d(2\ast2) + d(2+2)$ & $+(\alpha_2, \alpha_4)$ & $\alpha_5$ \\
$(2 \ast 2) \ast 2$ & $\ast(\alpha_1, \alpha_0)$ & $\alpha_6$ \\
$t$              & $+(\alpha_5, \alpha_6)$    & $\alpha_7$ \\
\end{tabular}
\end{center}

Note that some e-classes, e.g. $\alpha_0$ and $\alpha_1$, are shared between multiple e-nodes.

\paragraph{E-matching.}
We match the pattern $x \ast 2$---the left-hand side of the
equation---against the represented terms. There are \emph{two} matches, namely e-classes $\alpha_1$ and $\alpha_6$,
because two distinct e-classes are applications of $\ast$ whose right
argument is the class of $2$:
\[
\sigma_1 : x \mapsto \alpha_0, \qquad
\sigma_2 : x \mapsto \alpha_1.
\]
Each of these is really a valuation $a : 1 \to A_0 / R_0$, but at this point $R_0$ is the identity, so the valuations land in e-classes that are still singletons.

\paragraph{Adding $\hat\tau(t_2, \sigma)$.}
We collect both matches against the current e-graph, and only then apply them; this batching is what allows the invariants to be restored once, after all merges, rather than after each.
Firing the equation at a match $\sigma$ forms both instances and
identifies them. At $\sigma_1$ the instances are $\hat\tau(x\ast 2,
\sigma_1)$, the term $2\ast 2$ represented by $\ast(\alpha_0,\alpha_0)$ in class $\alpha_1$,
and $\hat\tau(x+x, \sigma_1)$, the term $2+2$ represented by $+(\alpha_0,\alpha_0)$ in class $\alpha_3$, already present;
firing merges their classes, calling $\union(\alpha_1, \alpha_3)$ so that $\alpha_1$ becomes the parent of $\alpha_3$.

At $\sigma_2$ the right-hand instance
$\hat\tau(x+x, \sigma_2)$ is the term $(2\ast2)+(2\ast2)$, represented by $+(\alpha_1, \alpha_1)$, which is \emph{not} in
$A_0$: firing \emph{adds} it as a new e-node, enlarging the represented set to
$A_1 \supsetneq A_0$, and invokes $\union(\alpha_6, \alpha_8)$, with $\alpha_8$ being the identifier for the new e-node.
Right after this step, and \emph{before} restoring the invariants, the affected
e-classes are as follows. The stored e-nodes still refer to their
original children---$\union$ has changed $U$ but not touched the
e-nodes---so $d(2{+}2)$ still reads $d(\alpha_3)$ even though
$\find(\alpha_3) = \alpha_1$:
\begin{center}
\begin{tabular}{c|c|c}
  Term & Stored e-node(s) & E-class\\
  \hline
  $\{\,2{+}2,\ 2{\ast}2\,\}$ & $\{\,{+}(\alpha_0, \alpha_0),\ {\ast}(\alpha_0, \alpha_0)\,\}$ & $\alpha_1$\\
  $d(2{\ast}2)$ & $d(\alpha_1)$ & $\alpha_2$\\
  $d(2{+}2)$ & $d(\alpha_3)$ & $\alpha_4$\\
  $\{\,(2{\ast}2){\ast}2,\ (2{\ast}2){+}(2{\ast}2)\,\}$ & $\{\,{\ast}(\alpha_1, \alpha_0),\ {+}(\alpha_1, \alpha_1)\,\}$ & $\alpha_6$\\
\end{tabular}
\end{center}
with the following parent links on identifiers
\[
\alpha_3 \xrightarrow{} \alpha_1 \qquad \alpha_8 \xrightarrow{} \alpha_6
\]
(read $b \to a$ as ``$a$ is the parent of $b$''), and every other
identifier being its own parent.
\paragraph{Restoring the invariants.}
We now need to make sure that the entries in the hashcons table are canonicalised e-nodes pointing to a canonical e-class identifier---this is necessary to maintain sharing---this is necessary to maintain sharing---this is necessary to maintain sharing---this is necessary to maintain sharing.
Furthermore, we need to make sure that the represented relation is a congruence, that is, if $\find(\alpha_i) = \find(\alpha_j)$ then $f(\alpha_i)$ and $f(\alpha_j)$ are in the same e-class.
Rebuilding re-canonicalises the children of each affected e-node through
$\find$ and re-inserts it. The node $d(\alpha_3)$ in class $\alpha_4$
canonicalises to $d(\alpha_1)$, since $\find(\alpha_3) = \alpha_1$; but
$d(\alpha_1)$ is already present in $H$, pointing to $\alpha_2$. This
collision is a violation of the congruence invariant---two applications
of $d$ to now-equal arguments---and is resolved by
$\union(\alpha_2, \alpha_4)$. The induced change to the child of
$+(\alpha_2, \alpha_4)$ in class $\alpha_5$ canonicalises it to
$+(\alpha_2, \alpha_2)$, with no further collision. After rebuild the
e-graph is:
\begin{center}
\begin{tabular}{c|c|c}
  Term & Canonical e-node(s) & E-class\\
  \hline
  $2$ & $2$ & $\alpha_0$\\
  $\{\,2{\ast}2,\ 2{+}2\,\}$ & $\{\,{\ast}(\alpha_0, \alpha_0),\ {+}(\alpha_0, \alpha_0)\,\}$ & $\alpha_1$\\
  $\{\,d(2{\ast}2),\ d(2{+}2)\,\}$ & $d(\alpha_1)$ & $\alpha_2$\\
  $d(2{\ast}2){+}d(2{+}2)$ & $+(\alpha_2, \alpha_2)$ & $\alpha_5$\\
  $\{\,(2{\ast}2){\ast}2,\ (2{\ast}2){+}(2{\ast}2)\,\}$ & $\{\,{\ast}(\alpha_1, \alpha_0),\ {+}(\alpha_1, \alpha_1)\,\}$ & $\alpha_6$\\
  $t$ & $+(\alpha_5, \alpha_6)$ & $\alpha_7$\\
\end{tabular}
\end{center}
The parent links are now
\[
\alpha_3 \xrightarrow{} \alpha_1, \qquad
\alpha_8 \xrightarrow{} \alpha_6, \qquad
\alpha_4 \xrightarrow{} \alpha_2,
\]
with every other identifier its own root.
Every entry of $H$ is now a
canonical e-node pointing to a canonical e-class identifier, and the
represented relation is a congruence.

We have seen the application of a single equation $(t_1, t_2)$ to an e-graph.
Because the equation is symmetric, we likewise match its right-hand side $x+x$ and merge those instances with the corresponding instances of $x \ast 2$.
Once an $A_i$ has been enlarged, or any merge performed, we re-scan all equations to see whether further matches have become available---a new term or a new merge may turn two previously distinct e-nodes into a redex.
This is repeated until a fixed point or a timeout is reached.
\end{example}

\begin{remark}
  In the above example, we assumed that the \emph{seed} of the e-graph is a single term.
  The process can be straightforwardly adjusted for the case with multiple terms in the seed.
  Another assumption that we made is that an e-graph operates with equations.
  In practice, e-graphs operate with rewrite rules, which are not necessarily symmetric, unlike equations.
  Our treatment results in a clear algebraic account of e-graphs, while rewrite rules do not necessarily define a \emph{theory} in the algebraic sense.
  We will maintain this assumption throughout the rest of the thesis.
\end{remark}

Finally, we have \emph{Lawvere theories}---introduced by
Lawvere~\cite{LawvereFunctorialSemanticsAlgebraic1963}---categories with
finite products generated by a single object. In the skeletal
presentation we use, the objects are the natural numbers, each $n$ being
the $n$-fold product of a generic object $1$. They provide a
presentation-independent view of algebraic theories: the same theory may
arise from different signatures and equations, but determines one Lawvere
theory. A morphism in $\mathcal{L}(n,1)$ is an $n$-ary operation, and a
morphism $n \to m$ an $m$-tuple of such operations; to capture a
\emph{theory} rather than a bare signature, these are taken up to
provable equality.

For $(\Sigma, \mathcal{E})$ the corresponding Lawvere theory
$\mathcal{L}_{(\Sigma,\mathcal{E})}$ is the finitely generated free
$(\Sigma,\mathcal{E})$-algebras organised into a category by
substitution. Its hom-set $\mathcal{L}_{(\Sigma,\mathcal{E})}(n,1) =
T_\Sigma(n)/\bar R$ is the carrier of the free
$(\Sigma,\mathcal{E})$-algebra on $n$ generators---the $\Sigma$-terms in
$n$ variables, which serve as the generators, modulo the congruence
$\bar R$ generated by $\mathcal{E}$---and a morphism $n \to m$ is an
$m$-tuple of these.

Composition is substitution. Given
\[
f = \langle f_1, \ldots, f_n \rangle \in \mathcal{L}_{(\Sigma,\mathcal{E})}(m,n),
\qquad
g = \langle g_1, \ldots, g_k \rangle \in \mathcal{L}_{(\Sigma,\mathcal{E})}(n,k),
\]
with $f_j \in T_\Sigma(m)/\bar R$ and $g_l \in T_\Sigma(n)/\bar R$, the
composite $g \circ f \in \mathcal{L}_{(\Sigma,\mathcal{E})}(m,k)$ is
obtained by substituting the $f_j$ for the variables of each $g_l$. This
substitution is the evaluation map in the free algebra
$T_\Sigma(m)/\bar R$,
\[
\hat\tau_n \;:\; \big(T_\Sigma(n)/\bar R\big) \times \big(T_\Sigma(m)/\bar R\big)^{n}
\longrightarrow T_\Sigma(m)/\bar R,
\]
giving
\[
g \circ f \;=\;
\big\langle\,
\hat\tau_n\big(g_1, (f_1, \ldots, f_n)\big),\ \ldots,\
\hat\tau_n\big(g_k, (f_1, \ldots, f_n)\big)
\,\big\rangle .
\]
Note that the components of $f$ and $g$ inhabit \emph{different} free
algebras, $T_\Sigma(m)/\bar R$ and $T_\Sigma(n)/\bar R$; composition is
therefore not the structure map of any single algebra, but the
substitution of $f$'s terms into $g$'s, carried out in
$T_\Sigma(m)/\bar R$.

The main benefit of this view is that it makes the notion of a model
parametric in the target category. The ordinary definition of a
$(\Sigma,\mathcal{E})$-algebra is tied to $\catname{Set}$: the carrier is
a set and the operations are functions on it. The Lawvere theory frees
the carrier from $\catname{Set}$: a model of the theory in \emph{any}
category $\mathcal{C}$ with finite products is a
finite-product-preserving functor $\mathcal{L}_{(\Sigma,\mathcal{E})} \to
\mathcal{C}$, recovering the ordinary $(\Sigma,\mathcal{E})$-algebras when
$\mathcal{C} = \catname{Set}$.

A canonical example is the theory of monoids interpreted in
$\catname{Cat}$. A finite-product-preserving functor $F :
\mathcal{L}_{\mathrm{Mon}} \to \catname{Cat}$ picks out a category
$\mathcal{A} = F(1)$, with $F(n) = \mathcal{A}^n$, and interprets the two
operations as
\[
\otimes := F(m) : \mathcal{A} \times \mathcal{A} \to \mathcal{A},
\qquad
F(e) : \mathbf{1} \to \mathcal{A},
\]
the latter a functor from the terminal category, i.e. a chosen object
$I \in \mathcal{A}$. The equations of the monoid theory, being equalities
of morphisms in $\mathcal{L}_{\mathrm{Mon}}$, are sent by $F$ to
\emph{equalities} of functors. The right unit law $m \circ (\id \times e)
= \id$ in $\mathcal{L}_{\mathrm{Mon}}(1,1)$, for instance, becomes
\[
\otimes \circ (\mathrm{id}_{\mathcal{A}} \times I) = \mathrm{id}_{\mathcal{A}},
\qquad\text{that is}\qquad
X \otimes I = X \ \text{ for all } X,
\]
and associativity $m \circ (m \times \id) = m \circ (\id \times m)$
becomes $(X \otimes Y) \otimes Z = X \otimes (Y \otimes Z)$. These hold
\emph{on the nose}, so $F$ is exactly a \emph{strict} monoidal category.

\begin{proposition}[Morphisms out of a free Lawvere theory]
\label{prop:lawvere-ump}
Let $\mathcal{L}_\Sigma$ be the free Lawvere theory on a signature
$\Sigma$, and let $\mathcal{L}$ be any Lawvere theory. Finite-product-preserving
functors $\mathcal{L}_\Sigma \to \mathcal{L}$ correspond naturally and
bijectively to interpretations of the generators: a choice, for each
operation $f \in \Sigma$ of arity $n$, of a morphism $\Phi(f) \in
\mathcal{L}(n,1)$. Each such choice extends uniquely, with $\Phi(n) =
\Phi(1)^n = n$ forced by product preservation and the action on all other
morphisms determined by functoriality.
\end{proposition}
This is the universal property of the free Lawvere theory: the functor
$F : \catname{Sig} \to \catname{Law}$ sending a signature to the theory
it freely generates is left adjoint to the underlying-signature functor
$U : \catname{Law} \to \catname{Sig}$, which sends a theory $\mathcal{L}$
to the signature whose operations of arity $n$ are the morphisms
$\mathcal{L}(n,1)$. The bijection
\[
\catname{Law}(\mathcal{L}_\Sigma,\ \mathcal{L}) \;\cong\;
\catname{Sig}(\Sigma,\ U\mathcal{L})
\]
is precisely the correspondence of the proposition.

A Lawvere theory also makes rewriting available on morphisms, and in a
form that absorbs a chosen equational modulus. Split the equations as
$\mathcal{E} = \mathcal{E}_0 \cup \mathcal{R}$, with $\mathcal{E}_0$ a
background theory and $\mathcal{R}$ oriented into rules. Working in
$\mathcal{L}_{(\Sigma,\mathcal{E}_0)}$, whose morphisms $n \to 1$ are
terms modulo $\mathcal{E}_0$, an oriented rule $\ell \to r$ rewrites these
morphisms by replacement under composition; the induced relation is
exactly rewriting modulo $\mathcal{E}_0$,
\[
f \to_{\mathcal{R}/\mathcal{E}_0} g
\iff
\exists\, h, k:\quad f =_{\mathcal{E}_0} h,\quad h \to_{\mathcal{R}} k,\quad k =_{\mathcal{E}_0} g,
\]
since $=_{\mathcal{E}_0}$ is precisely equality of morphisms in
$\mathcal{L}_{(\Sigma,\mathcal{E}_0)}$. This is the same system as
rewriting the terms of the free $(\Sigma,\mathcal{E}_0)$-algebra
$T_\Sigma(n)/\bar R_{\mathcal{E}_0}$ by $\mathcal{R}$, the morphisms
$n \to 1$ being exactly those terms.
% This is yet another way of looking at e-graphs: picking some equations $\mathcal{E}_0$ and some rules $\mathcal{R}$ that become the rewrite rules, an e-graph performs rewritings steps from $\mathcal{R}$ modulo $\mathcal{E}_0$.

\begin{remark}
Throughput this chapter we have considered only \emph{single-sorted} algebraic theories, with a single carrier set meaning that every variable has the same type.
The generalisation to \emph{multi-sorted} algebraic theories is straightforward: the signature $\Sigma$ is taken to be $\Sigma : \mathbb{S}^{*} \times \mathbb{S} \to \catname{Set}$ with the idea that each $f \in \Sigma([s_1, \ldots, s_m], s)$ is an operation with $m$ inpputs typed $s_1, \ldots, s_m$ and output type $s$.
The carrier of a corresponding $\Sigma$-algebra is a family of sets indexed by $\mathbb{S}$.
The rest is built up on that analogously to a single-sorted case.
\end{remark}